% Part: sets-functions-relations
% Chapter: relations
% Section: operations

\documentclass[../../../include/open-logic-section]{subfiles}

\begin{document}

\olfileid[ar]{sfr}{rel}{ops}
\olsection{عمليات على العلاقات}

قد ينفع أن نغير العلاقات أو نضم بعضها إلى بعض. ففي
\olref[sfr][rel][ord]{prop:stricttopartial} أخذنا \emph{اتحاد} علاقتين،
وهو اتحادهما حين تؤخذان مجموعتين من الأزواج. وفي
\olref[sfr][rel][ord]{prop:partialtostrict} أخذنا الفرق النسبي
بينهما. وها هنا عمليات أخرى على العلاقات.

\begin{defn}\ollabel{relationoperations}
خذ علاقتين $R$ و$S$، ومجموعة كيفما كانت $A$.

نسمّي \emph{معكوس} $R$ المجموعة $R^{-1} = \Setabs{\tuple{y, x}}{\tuple{x,
    y} \in R}$.

ونسمّي \emph{الضرب النسبي} لـ$R$ و$S$ المجموعة $(R \mid S) =
\{\tuple{x, z} : \exists y(Rxy \land Syz)\}$.

و\emph{تقييد} $R$ على $A$ معرّف بـ$\funrestrictionto{R}{A}= R
\cap A^2$.

و\emph{تطبيق} $R$ على $A$ معرّف بـ$\funimage{R}{A} = \{y :
(\exists x \in A)Rxy\}$.
\end{defn}

\begin{ex}
خذ علاقة اللاحق $S \subseteq \Int^2$ على~$\Int$، وهي
$S = \Setabs{\tuple{x, y} \in \Int^2}{x + 1 = y}$؛ فـ$Sxy$ يكافئ $x + 1 = y$.

فمعكوسها $S^{-1}$ علاقة السابق على $\Int$، أي
$\Setabs{\tuple{x,y}\in\Int^2}{x -1 =y}$.

وأما $S\mid S$ فتعطي
$ \Setabs{\tuple{x,y}\in\Int^2}{x + 2 =y}$.

وتقييدها $\funrestrictionto{S}{\Nat}$ علاقة اللاحق على~$\Nat$.

وتطبيقها يعطي $\funimage{S}{\{1,2,3\}}$ مساوية لـ$\{2, 3, 4\}$.
\end{ex}

\begin{defn}[الإغلاق المتعدي]
خذ العلاقة الثنائية $R \subseteq A^2$.

نعرّف \emph{الإغلاق المتعدي} لـ~$R$ بأنه $R^+ = \bigcup_{0 < n \in
\Nat} R^n$؛ والقوى فيه معرفة عوديًا: $R^1 = R$ و$R^{n+1} = R^n
\mid R$.

وأما \emph{الإغلاق الانعكاسي المتعدي} لـ$R$ فتعريفه $R^* = R^+ \cup
\Id{A}$.
\end{defn}

\begin{ex}
في علاقة اللاحق $S \subseteq \Int^2$ يكون $S^2xy$ مكافئًا لـ$x + 2 =
y$، و$S^3xy$ مكافئًا لـ$x + 3 = y$، وعلى هذا القياس.
فـ$S^+xy$ إذا وفقط إذا كان $x + n = y$ لعدد $n \geq 1$.
وحاصل ذلك أن $S^+xy$ يكافئ $x < y$، وأن $S^*xy$ يكافئ $x \le
y$.
\end{ex}

\begin{prob}
بيّن أن إغلاق $R$ المتعدي يحقق التعدية حقًا.
\end{prob}

\end{document}
